% ============================================================
% Section 2: Related Work
% ============================================================

\section{Related Work}
\label{sec:related}

\subsection{Document and Image Forgery Detection}

Early work on image forensics relied on hand-crafted features such as noise
inconsistencies, JPEG compression artifacts, and chromatic aberration patterns to
detect copy-move and splicing manipulations~\cite{huh2018fighting}. Deep learning
substantially shifted the field: convolutional neural networks can learn to detect
subtle statistical fingerprints left by image-processing operations, often
outperforming engineered pipelines on standard benchmarks~\cite{bappy2019hybrid,
mantranet2019}.

Recent large-scale foundations have further advanced the state of the art.
\textsc{TruFor}~\cite{trufor2023} pre-trains on over 876\,000 images and learns to
extract high-level semantic inconsistencies as well as low-level noise-pattern
anomalies, achieving strong performance on general image forensics benchmarks.
However, such general-purpose models are optimized for natural image manipulation
and do not account for the structured layout, typography, and holographic elements
specific to identity documents. Applying them to the identity-document domain
requires fine-tuning and does not exploit document-specific priors.

Identity document forgery introduces unique challenges compared to natural-image
manipulation: (i)~the forged regions are often small (a single text field or a
face) relative to the full document, (ii)~the background texture is highly
structured (guilloch\'e patterns, UV-reactive ink), which creates strong false
negatives for general-purpose noise detectors, and (iii)~the diversity of
document types across countries and issuers demands generalization beyond what
natural-image pre-training provides. The DeepID 2025 Challenge~\cite{deepid2025challenge}
was specifically designed to benchmark progress on this task using the
FantasyID dataset~\cite{fantasyid2024}, which contains synthetic attack types
including face swaps and text inpainting on fantasy identity documents.

\subsection{Multi-Task Learning for Detection and Localization}

Joint detection and localization is a canonical instance of multi-task learning
(MTL)~\cite{caruana1997multitask, ruder2017overview}. In MTL, shared representations
learned for one task regularize learning for related tasks, typically yielding better
generalization than independent training when tasks are correlated~\cite{crawshaw2020multi}.

The dominant approach to combining task-specific losses is scalar
scalarization~\cite{marler2004survey}: a single objective $\mathcal{L} = \sum_t
\lambda_t \mathcal{L}_t$ is minimized, with fixed weights $\lambda_t$ set before
training. Kendall \textit{et al.}~\cite{kendall2018multi} proposed learning $\lambda_t$
jointly with the model weights via homoscedastic uncertainty, but this requires an
additional learning-rate schedule and is sensitive to initialization. Liu \textit{et al.}~\cite{liu2019end}
introduced gradient-normalization to dynamically balance task gradients, at the cost
of additional computation per step. All scalar approaches share a common limitation:
the optimal trade-off between tasks is not established before training, yet the
weights implicitly commit to a fixed trade-off.

In the image forensics domain, U-Net~\cite{unet2015} architectures have been widely
adopted for pixel-level localization because their skip connections preserve spatial
resolution while the encoder captures high-level semantic anomaly features. Several
works combine classification and segmentation branches in a single network, but
typically fix the task weights empirically without systematic
optimization~\cite{wu2022robust}. Our work differs by including the task weight
$\lambda_{\text{mask}}$ in the HPO search space and selecting it via a
multi-objective criterion rather than a grid search.

\subsection{Hyperparameter Optimization and Multi-Objective Search}

Hyperparameter optimization (HPO) has evolved from manual search and grid search to
sequential model-based methods. Tree-structured Parzen Estimator
(TPE)~\cite{tpe2011} models the distribution of good and bad configurations
separately and selects candidates that maximize the ratio of the two models, a
strategy significantly more sample-efficient than random search~\cite{randomsearch2012}.
The Optuna framework~\cite{optuna2019} implements TPE as its default sampler and
has become a standard tool for HPO in deep learning.

Multi-objective HPO extends single-objective search to the case where multiple
metrics must be jointly optimized. The output is a Pareto-optimal set of
configurations rather than a single best point, making the trade-off between
objectives explicit~\cite{miettinen1999nonlinear}. Evolutionary algorithms such as
NSGA-II~\cite{nsga2_2002} are the most established approach, but they require
large population sizes to maintain Pareto-front diversity. The Multiobjective TPE
(MOTPE) sampler~\cite{motpe2022} adapts the TPE surrogate model to the
multi-objective setting, achieving competitive Pareto-front coverage with far fewer
function evaluations than evolutionary methods—a critical advantage when each
evaluation requires training a deep neural network.

To our knowledge, this is the first work to apply Pareto-based multi-objective HPO
to the problem of joint identity-document forgery detection and localization, and
the first to use the resulting Pareto front as the primary model-selection criterion
in a nested cross-validation framework.
